\section{Ecuaciones diferenciales, solución y aproximación}
Una ecuación diferencial parcial (PDE, por sus siglas en inglés) describe relaciones locales entre una magnitud y sus variaciones. Temperatura, velocidad, desplazamiento y potencial eléctrico son ejemplos de campos: su valor cambia con la posición y, en algunos problemas, con el tiempo. La ecuación por sí sola no determina la respuesta física. También intervienen el dominio, los parámetros materiales, las fuentes y las condiciones de frontera e iniciales \cite{evans1998pde}.

En problemas particulares existen soluciones analíticas que permiten expresar el campo mediante funciones conocidas. En otros se recurre a discretizaciones como diferencias, volúmenes o elementos finitos, o a aproximadores neuronales sujetos a las ecuaciones. La elección depende de geometría, regularidad, precisión requerida y recursos; una representación neuronal no reemplaza automáticamente los métodos clásicos \cite{fenics2016,grossmann2024}.

Sean $\Omega\subset\mathbb R^d$ el dominio espacial, $\partial\Omega$ su frontera y $\mathbf u$ un vector de $m$ campos. Escribiendo $z=x$ para un problema estacionario y $z=(x,t)$ para uno evolutivo, un problema general tiene la forma
\begin{equation}
\mathcal N_a[\mathbf u;\mu](z)=f_a(z),\qquad
\mathcal B_b[\mathbf u](z)=g_b(z),\qquad
\mathbf u(x,0)=\mathbf u_0(x).
\label{eq:pde-general}
\end{equation}
El índice $a$ recorre ecuaciones y $b$ condiciones; no necesariamente hay una condición por cada ecuación. $\mu$ contiene parámetros físicos. Dirichlet prescribe valores; Neumann prescribe derivadas normales o flujos, según el operador; Robin combina ambos. En elasticidad, la condición natural es una tracción, no una derivada de la tensión. El problema estacionario omite la condición inicial.

\subsection{Familias de operadores y escalas de la solución}
La ecuación de Poisson, $-\Delta u=f$, es el prototipo elíptico: relaciona curvatura espacial y fuente en un estado de equilibrio. Una ecuación de difusión $u_t-d\Delta u=f$, con $d>0$, incorpora evolución y tiende a suavizar variaciones espaciales. En una ley de transporte, $u_t+\mathbf a\cdot\nabla u=0$, la información se desplaza con velocidad $\mathbf a$. Cuando el transporte depende de la propia solución, como en $u_t+u u_x=\nu u_{xx}$, puede competir con la difusión y formar transiciones estrechas. Las clasificaciones elíptica, parabólica e hiperbólica describen propiedades del operador principal; no son etiquetas de dificultad computacional universal \cite{evans1998pde}.

Dos soluciones igualmente suaves pueden requerir resoluciones diferentes. Para $u(x)=\sin(k\pi x)$, derivar multiplica la amplitud por $k\pi$ y derivar dos veces por $(k\pi)^2$. Una capa de anchura $\ell$ produce derivadas que escalan aproximadamente como $\ell^{-1}$ y $\ell^{-2}$. La dificultad no reside únicamente en el tamaño del campo: una amplitud pequeña puede coexistir con gradientes grandes. En sistemas acoplados, además, un campo puede tener una escala característica distinta de los demás.

\subsection{Error, residual y estabilidad}
Para una aproximación $\mathbf u_\theta$, el error y el residual son objetos diferentes:
\begin{equation}
\mathbf e=\mathbf u_\theta-\mathbf u,\qquad
r_a=\mathcal N_a[\mathbf u_\theta;\mu]-f_a.
\label{eq:error-residual}
\end{equation}
El error requiere conocer la solución o una referencia; el residual puede evaluarse con el modelo y la ecuación. Un residual pequeño en ciertos puntos no garantiza error pequeño en todo el dominio. Una frontera mal satisfecha, una capa no muestreada o un problema mal condicionado pueden invalidar esa interpretación \cite{deryckmishra2022,grossmann2024}.

Las normas precisan qué significa cercanía. Para un campo escalar,
\begin{equation}
\|e\|_{L^2(\Omega)}^2=\int_\Omega |e|^2dx,\qquad
|e|_{H^1(\Omega)}^2=\int_\Omega\|\nabla e\|^2dx.
\end{equation}
La primera agrega errores de valor; la segunda detecta errores de gradiente. $H^1$ contiene funciones cuadrado integrables con primeras derivadas débiles cuadrado integrables. La derivada débil extiende la integración por partes a funciones que no poseen derivadas clásicas en todos los puntos. Esta distinción permite representar soluciones con cambios de material sin exigir suavidad artificial a través de la interfaz.

Para Poisson con frontera homogénea, $-\Delta e=r$ en sentido débil. Probar con $e$ produce $\|\nabla e\|^2=\langle r,e\rangle$. Por la norma dual,
\begin{equation}
\|\nabla e\|_{L^2}\leq\|r\|_{H^{-1}},
\end{equation}
si $H_0^1$ se dota de la norma del gradiente. Aquí sí existe una relación residual--error, pero utiliza una norma funcional, toda la ecuación y una frontera correcta. Una media empírica de $r^2$ no es esa norma dual. Para otros operadores cambian las constantes, normas e hipótesis de estabilidad.

\subsection{Adimensionalización y significado de las pérdidas}
Una suma de errores en metros, voltios y pascales depende de la elección de unidades. Sean $x=L\hat x$, $t=T\hat t$ y $u=U\hat u$. La regla de la cadena da
\begin{equation}
u_t=\frac UT\hat u_{\hat t},\qquad
u_x=\frac UL\hat u_{\hat x},\qquad
u_{xx}=\frac U{L^2}\hat u_{\hat x\hat x}.
\label{eq:escalas}
\end{equation}
En Burgers resulta $\hat u_{\hat t}+(UT/L)\hat u\hat u_{\hat x}=(\nu T/L^2)\hat u_{\hat x\hat x}$. Elegir $T=L/U$ deja el grupo $\nu/(UL)$, inverso de un número de Reynolds. Adimensionalizar no elimina la competencia física; la expresa mediante cocientes que conservan significado al cambiar unidades. Normalizar coordenadas y normalizar residuales son operaciones relacionadas, pero diferentes \cite{wang2023expertguide}.

% Fragmento listo para integrarse en el marco teórico.
% Las cuatro figuras se esperan en la ruta de figuras del documento principal:
% p1-poisson.pdf, p2-burgers.pdf, p3-allen-cahn.pdf y p4-dfg.pdf.
% Este fragmento interpreta los contratos de P1--P4; no sustituye la
% especificación experimental del capítulo de metodología.

\section{Lectura física de los problemas canónicos}
\label{sec:lectura-fisica-pdes}

Las ecuaciones diferenciales parciales del proyecto no son ejemplos
intercambiables. Cada una concentra una dificultad física distinta: Poisson
representa un equilibrio elíptico con variación distribuida; Burgers separa
transporte no lineal y difusión en una transición estrecha; Allen--Cahn
combina reacción, difusión y capas de interfaz; y el benchmark DFG exige
respetar simultáneamente inercia, viscosidad, incompresibilidad, geometría y
balances integrales. Esta lectura sirve para entender qué significa que una
aproximación represente bien un campo. Los contratos exactos del banco de
problemas se conservan en el capítulo de metodología; aquí se explican sus
consecuencias físicas y las magnitudes que hacen visibles.

Las figuras de esta sección son esquemas matemáticos construidos a partir de
las ecuaciones y parámetros declarados. Cuando se muestran mapas o curvas
obtenidos al evaluar una fórmula exacta, se trata de una ilustración analítica
del problema y no de un resultado de entrenamiento. En particular, ninguna
figura de esta sección certifica por sí sola una mejora de una red o una
convergencia numérica.

\subsection{P1. Poisson: equilibrio y frecuencia distribuida}

Considérese el cuadrado unitario $\Omega=(0,1)^2$ y el problema
\begin{equation}
-\Delta u=f_k \quad \text{en }\Omega,
\qquad u=0 \quad \text{en }\partial\Omega,
\qquad
u_k^*(x,y)=\sin(k\pi x)\sin(k\pi y),
\label{eq:expansion-p1-poisson}
\end{equation}
donde
\begin{equation}
f_k(x,y)=2(k\pi)^2u_k^*(x,y).
\label{eq:expansion-p1-source}
\end{equation}
La función $u_k^*$ se anula en los cuatro lados, por lo que satisface la
condición de Dirichlet sin corrección adicional. Además,
$\partial_{xx}u_k^*=-(k\pi)^2u_k^*$ y
$\partial_{yy}u_k^*=-(k\pi)^2u_k^*$; por tanto,
$-\Delta u_k^*=2(k\pi)^2u_k^*=f_k$. La igualdad es una verificación por
sustitución de la solución manufacturada, no una aproximación de los valores
de $u$.

La interpretación física más simple es un estado estacionario: el operador
$-\Delta$ mide cómo la curvatura del campo compensa una fuente distribuida.
Si $u$ representa temperatura en un medio homogéneo, $f$ puede representar
generación interna de calor; si representa un potencial, $f$ desempeña el
papel de una densidad de fuente. En ambos casos no hay una dirección de
propagación ni una condición inicial. La frontera fija el nivel del campo y
el interior ajusta su curvatura para satisfacer el balance.

El parámetro $k$ cambia la escala espacial sin cambiar la amplitud máxima de
$u_k^*$. Los ceros consecutivos de $\sin(k\pi x)$ están separados por
$1/k$, mientras que la curvatura y la fuente crecen como $k^2$. Para leer
la oscilación sin mezclarla con una región de interés arbitraria, puede
tomarse el corte
\begin{equation}
u_k^*\!\left(x,\frac{1}{2k}\right)=\sin(k\pi x),
\label{eq:expansion-p1-corte}
\end{equation}
porque $\sin(k\pi/(2k))=1$. Así, $k=1$ muestra una variación amplia y
$k=4$ muestra cuatro semiondas en la misma longitud. La familia permite
preguntar si la representación resuelve frecuencia distribuida en todo el
dominio; una ganancia en P1 no puede interpretarse automáticamente como
localización de capacidad en una interfaz o frente.

\begin{figure}[htbp]
\centering
\includegraphics[width=\linewidth]{p1-poisson.pdf}
\caption[Poisson: frecuencia distribuida y corte exacto.]{Esquema analítico
de P1 construido con $u_k^*(x,y)=\sin(k\pi x)\sin(k\pi y)$ para $k=1$ y
$k=4$. Los mapas muestran la variación en $\Omega=(0,1)^2$; los cortes se
evalúan en $y=1/(2k)$, donde la ordenada coincide exactamente con
$\sin(k\pi x)$. La figura es una ilustración de la solución manufacturada, no
un resultado de entrenamiento ni una medición FEM.}
\label{fig:expansion-p1-poisson}
\end{figure}

La forma débil hace explícita la misma lectura. Para una prueba $v$ que se
anula en $\partial\Omega$,
\begin{equation}
\int_\Omega \nabla u\cdot\nabla v\,\mathrm d x
 =\int_\Omega f_kv\,\mathrm d x.
\label{eq:expansion-p1-debil}
\end{equation}
El lado izquierdo mide el trabajo asociado a los gradientes y el derecho la
acción de la fuente. Una discretización puede aproximar esta identidad aun
cuando el campo local tenga error; por ello se deben distinguir balance
integral, error de campo y error de derivadas. Esta distinción es central al
comparar una aproximación neuronal con una referencia clásica
\cite{evans1998pde,fenics2016}.

\subsection{P2. Burgers: transporte no lineal, difusión y frente}

En $x\in[-1,1]$ y $t\in[0,1]$, el segundo contrato es
\begin{equation}
u_t+u\,u_x-\nu u_{xx}=0,
\qquad
u^*(x,t)=c-A\tanh q,
\qquad
q=\frac{A(x-ct-x_0)}{2\nu}.
\label{eq:expansion-p2-burgers}
\end{equation}
La condición inicial y las dos condiciones de frontera son las trazas de
$u^*$:
\begin{equation}
u(x,0)=u^*(x,0),
\qquad
u(-1,t)=u^*(-1,t),
\qquad
u(1,t)=u^*(1,t).
\label{eq:expansion-p2-conditions}
\end{equation}
Se fijan $A=0.5$, $x_0=-0.25$,
$c\in\{0,0.5\}$ y $\nu\in\{0.02,0.005\}$. El centro del frente está en
$x_f(t)=ct+x_0$: con $c=0$ permanece estacionario y con $c=0.5$ recorre
media unidad durante el intervalo temporal. La frontera dependiente del
tiempo no es un forzamiento añadido; es la traza necesaria para mantener la
onda exacta al entrar o salir del dominio.

La sustitución muestra qué términos se equilibran. Si
$S=\operatorname{sech}^2q$, entonces
\begin{equation}
u_x=-\frac{A^2}{2\nu}S,
\qquad
u_t=-c\,u_x,
\qquad
u_{xx}=\frac{A^3}{2\nu^2}S\tanh q.
\label{eq:expansion-p2-derivatives}
\end{equation}
Como $u-c=-A\tanh q$,
\begin{equation}
u_t+u u_x=(u-c)u_x
 =\frac{A^3}{2\nu}S\tanh q
 =\nu u_{xx}.
\label{eq:expansion-p2-balance}
\end{equation}
El término $u u_x$ transporta la forma mediante la velocidad local; la
difusión $\nu u_{xx}$ impide que la transición se vuelva discontinua. El
ancho característico $\ell=2\nu/A$ vale $0.08$ cuando $\nu=0.02$ y
$0.02$ cuando $\nu=0.005$. El máximo del gradiente tiene magnitud
$A^2/(2\nu)$, que pasa de $6.25$ a $25$ al reducir $\nu$. La amplitud del
campo apenas cambia, pero la derivada que entra en el residual se concentra
en una región cuatro veces más estrecha.

Esta relación explica por qué P2 es una prueba de resolución localizada. Una
aproximación puede tener valores razonables fuera del frente y, al mismo
tiempo, perder la curvatura que equilibra transporte y difusión. También
explica por qué movimiento y estrechamiento son dificultades distintas: $c$
cambia la posición de la señal en el espacio-tiempo, mientras $\nu$ cambia
la escala de sus derivadas. La familia y sus ondas viajeras tienen antecedente
analítico en la literatura de Burgers \cite{burgerswaves2025}; los cuatro
cruces de parámetros pertenecen al contrato finito de este proyecto.

\begin{figure}[htbp]
\centering
\includegraphics[width=\linewidth]{p2-burgers.pdf}
\caption[Burgers: frente, movimiento y balance de términos.]{Diagnóstico
analítico de P2 evaluado en la onda exacta. Se comparan los perfiles iniciales
para $\nu=0.02$ y $\nu=0.005$ con $c=0.5$, se muestra la concentración de
$|u_x|$, se sigue la trayectoria $x_f(t)=ct+x_0$ para $\nu=0.02$ y se
comparan $u_t$, $u u_x$ y $-\nu u_{xx}$ en $t=0.5$. Las curvas satisfacen
la identidad de balance por construcción; la figura es una ilustración de la
solución exacta, no una corrida de PINN o PIKAN.}
\label{fig:expansion-p2-burgers}
\end{figure}

\subsection{P3. Allen--Cahn: reacción, energía e interfaz}

El tercer problema conserva el contrato del ejemplo documentado de DeepXDE:
\begin{align}
u_t-d u_{xx}-5(u-u^3)&=0,
&&x\in[-1,1],\quad t\in[0,1],\quad d=0.001,\\
u(x,0)&=x^2\cos(\pi x),
&&u(-1,t)=u(1,t)=-1.
\label{eq:expansion-p3-allen-cahn}
\end{align}
Puede escribirse como un flujo de gradiente de la energía
\begin{equation}
\mathcal E[u]=\int_{-1}^{1}
\left[\frac d2|u_x|^2+W(u)\right]\,\mathrm d x,
\qquad
W(u)=\frac54(u^2-1)^2.
\label{eq:expansion-p3-energy}
\end{equation}
En efecto, $W'(u)=5(u^3-u)$ y la ecuación equivale a
$u_t=d u_{xx}-W'(u)$. La reacción $5(u-u^3)$ empuja el campo hacia los
estados estables $u=\pm1$ y hace inestable el estado $u=0$; el término
difusivo penaliza cambios espaciales demasiado bruscos. La solución debe
resolver ambos impulsos, además de acomodarse al dato inicial y a las
fronteras fijadas en $-1$.

La escala de una interfaz se obtiene al equilibrar difusión y reacción. Para
el perfil estacionario idealizado que conecta $-1$ con $+1$,
\begin{equation}
u_{\rm kink}(x)=
\tanh\!\left(\frac{x-x_c}{\ell_{\rm AC}}\right),
\qquad
\ell_{\rm AC}=\sqrt{\frac{2d}{5}}=0.02.
\label{eq:expansion-p3-kink}
\end{equation}
La cifra $0.02$ es la escala del perfil de transición para $d=0.001$; no es
una afirmación de que la evolución con el dato inicial del contrato tenga una
única interfaz de ese ancho en todo tiempo. El estado inicial
$x^2\cos(\pi x)$ puede producir varias zonas de transición, y las
condiciones $u(\pm1,t)=-1$ influyen en su evolución. La máscara física de
una interfaz debe construirse a partir del campo de referencia cuando se
evalúa el problema, no a partir de la forma idealizada de esta explicación.

\begin{figure}[htbp]
\centering
\includegraphics[width=\linewidth]{p3-allen-cahn.pdf}
\caption[Allen--Cahn: potencial, reacción y perfil de interfaz.]{Esquema
analítico de P3 para $d=0.001$: el potencial $W(u)=\tfrac54(u^2-1)^2$
muestra sus mínimos en $u=\pm1$, la reacción $5(u-u^3)=-W'(u)$ muestra
los estados favorecidos y el perfil $\tanh((x-x_c)/\sqrt{2d/5})$ ilustra
la escala de un frente estacionario. No se muestra una evolución numérica ni
se atribuye al panel un error de referencia; el dato inicial y las fronteras
del contrato son los de \eqref{eq:expansion-p3-allen-cahn}.}
\label{fig:expansion-p3-allen-cahn}
\end{figure}

Allen--Cahn introduce una diferencia importante respecto de Burgers. En
Burgers la transición se mueve por un balance entre advección y difusión; en
Allen--Cahn la transición también puede cambiar porque el campo desciende por
un potencial local con dos estados preferidos. Por eso el valor de $u$ y su
gradiente tienen interpretaciones simultáneas: estar cerca de $\pm1$ indica
fase, mientras que $|u_x|$ identifica una interfaz. Una representación que
suaviza demasiado la interfaz puede reducir una medida global y aun así
desplazar el lugar donde el campo cambia de fase. La formulación y las
condiciones se mantienen coherentes con el ejemplo de DeepXDE versión 1.12.2
\cite{deepxdeAC}.

\subsection{P4. DFG 2D-1: balance de cantidad de movimiento}

El benchmark DFG 2D-1 considera flujo estacionario incompresible alrededor de
un cilindro. El dominio es el canal
\begin{equation}
\Omega=([0,2.2]\times[0,0.41])
\setminus\overline{B_{0.05}(0.2,0.2)},
\label{eq:expansion-p4-domain}
\end{equation}
con densidad unitaria, viscosidad cinemática $\nu=0.001$ y campos
$\mathbf u=(u,v)$ y $p$. El sistema local es
\begin{align}
u u_x+v u_y+p_x-\nu(u_{xx}+u_{yy})&=0,\\
u v_x+v v_y+p_y-\nu(v_{xx}+v_{yy})&=0,\\
u_x+v_y&=0.
\label{eq:expansion-p4-ns}
\end{align}
Los dos primeros renglones expresan el balance de momento: la convección
transporta momento, el gradiente de presión lo redistribuye y la viscosidad
disipa gradientes. El tercero impone que el volumen no se comprima; no es una
ecuación auxiliar que pueda omitirse sin cambiar el problema.

La entrada del canal prescribe el perfil parabólico
\begin{equation}
\mathbf u_{\rm in}(y)=
\left(\frac{4(0.3)y(0.41-y)}{0.41^2},0\right),
\label{eq:expansion-p4-inlet}
\end{equation}
cuya velocidad máxima es $0.3$ en la línea media y cuya media transversal
es $0.2$. Las paredes del canal y la circunferencia del cilindro tienen
condición de no deslizamiento, $\mathbf u=\mathbf0$. En la salida se usa,
con la convención del benchmark para el operador Laplaciano,
\begin{equation}
\nu\,\partial_n\mathbf u-p\mathbf n=\mathbf0.
\label{eq:expansion-p4-outlet}
\end{equation}
Esta condición deja salir el flujo sin fijar artificialmente un valor de
presión adicional. Con $L=0.1$ como diámetro y $U=0.2$ como velocidad de
referencia, $Re=UL/\nu=20$, de modo que el caso es laminar y permite
examinar la estela sin convertir el problema en un estudio de turbulencia.

La forma integral conecta la PDE con las cantidades de ingeniería. Como
$\nabla\cdot\mathbf u=0$ y $\nu$ es constante, para un volumen de control
$K\subset\Omega$ sin fuerza volumétrica se obtiene
\begin{equation}
\int_{\partial K}
\left[\mathbf u(\mathbf u\cdot\mathbf n)+p\mathbf n
      -\nu(\nabla\mathbf u)\mathbf n\right]\,\mathrm d s=\mathbf0.
\label{eq:expansion-p4-integral}
\end{equation}
Cada cara intercambia flujo convectivo, presión y esfuerzo viscoso. Sobre el
cilindro, la tracción según la convención usada en el contrato es
$\boldsymbol\sigma_b\mathbf n$, con
$\boldsymbol\sigma_b=\nu\nabla\mathbf u-pI$. Por tanto,
\begin{equation}
\begin{pmatrix}F_D\\F_L\end{pmatrix}
 =\int_S\boldsymbol\sigma_b\mathbf n\,\mathrm d s,
\qquad
C_D=\frac{2F_D}{U^2L},
\qquad
C_L=\frac{2F_L}{U^2L}.
\label{eq:expansion-p4-forces}
\end{equation}
El benchmark publica valores de arrastre, sustentación y caída de presión que
se usan como controles externos en la metodología; no son resultados
producidos por esta propuesta \cite{dfgRe20}.

\begin{figure}[htbp]
\centering
\includegraphics[width=\linewidth]{p4-dfg.pdf}
\caption[DFG 2D-1: geometría, entrada y balance.]{Esquema matemático del
benchmark DFG 2D-1 a $Re=20$. Se identifica el canal
$[0,2.2]\times[0,0.41]$, el cilindro de radio $0.05$ centrado en
$(0.2,0.2)$, las paredes, la entrada y la salida; se añade el perfil
parabólico de entrada y un panel de los términos convectivo, de presión y
viscoso del balance. La figura no contiene líneas de corriente ni una
solución calculada: su función es fijar geometría, condiciones y signos del
contrato.}
\label{fig:expansion-p4-dfg}
\end{figure}

El carácter canónico de DFG proviene de que un campo visualmente plausible no
garantiza el balance. Una pequeña violación de $u_x+v_y=0$ acumula flujo
espurio; una presión con nivel o signo inconsistente altera la tracción; y
una salida mal impuesta puede contaminar la estela. Por eso el desempeño debe
leerse en campos y en observables integrales: arrastre, sustentación y caída
de presión. La geometría curva y el no deslizamiento junto al cilindro
introducen gradientes donde una descripción global suave puede parecer
adecuada lejos de la superficie y fallar en la región que determina las
fuerzas.

\subsection{Lectura conjunta del cuarteto}

P1 cambia la frecuencia en todo el dominio; P2 concentra una derivada en un
frente que puede moverse; P3 combina un potencial local con una interfaz cuya
escala se obtiene de un balance reacción--difusión; y P4 exige compatibilidad
entre campos, frontera curva y cantidades integrales. La secuencia permite
separar cuatro preguntas físicas: si una representación resuelve variación
distribuida, si conserva una transición transportada, si mantiene estados y
interfaces, y si satisface un balance vectorial con incompresibilidad.

Las cuatro figuras acompañan esas preguntas con objetos distintos: funciones
manufacturadas en P1, derivadas exactas en P2, energía y perfil estacionario
en P3, y geometría más balance en P4. Su propósito es preparar la lectura de
las métricas y referencias del capítulo de metodología. La evidencia de una
hipótesis sobre adaptación solo podrá provenir de comparaciones reproducibles
con esas referencias; las ilustraciones presentes aquí explican el problema,
pero no sustituyen tal evidencia.


\section{Métodos clásicos y adaptación de la resolución}
Una discretización reemplaza un espacio infinito de funciones por un conjunto finito de incógnitas. La malla organiza posiciones, celdas o elementos; las funciones de aproximación determinan qué variaciones pueden representarse. Diferencias finitas (FDM), volúmenes finitos (FVM) y elementos finitos (FEM), por sus siglas en inglés, proporcionan tanto soluciones numéricas como antecedentes conceptuales de la adaptación \cite{fenics2016}.

\subsection{Diferencias finitas: de Taylor al operador discreto}
Para una función suficientemente suave,
\begin{align}
u(x+h)&=u+hu'+\tfrac{h^2}{2}u''+\tfrac{h^3}{6}u'''+O(h^4),\\
u(x-h)&=u-hu'+\tfrac{h^2}{2}u''-\tfrac{h^3}{6}u'''+O(h^4).
\end{align}
Al sumar, las potencias impares se cancelan. Despejar $u''$ produce
\begin{equation}
u''(x_i)=\frac{u_{i-1}-2u_i+u_{i+1}}{h^2}+O(h^2).
\end{equation}
El término $O(h^2)$ expresa consistencia local bajo suavidad, no una garantía aislada de convergencia del problema completo. Fronteras, estabilidad del integrador temporal y resolución de capas también importan. Reducir la tolerancia de un solucionador algebraico no corrige por sí solo el error de discretización.

\subsection{Volúmenes finitos y conservación}
Integrar una ley $u_t+\nabla\cdot F(u)=s$ en una celda $K$ y aplicar el teorema de la divergencia da
\begin{equation}
\frac{d}{dt}\int_Ku\,dx+\int_{\partial K}F(u)\cdot n_K\,ds=\int_Ks\,dx.
\end{equation}
El método sustituye el flujo de frontera por un flujo numérico compartido entre celdas. Las contribuciones internas se cancelan al sumar balances, siempre que se utilice el mismo flujo con normales opuestas. Esta propiedad explica su interés para transporte y fluidos. Conservar un balance integral no equivale a obtener un campo local exacto: ambas propiedades deben distinguirse.

\subsection{Elementos finitos: forma débil, base y ensamblaje}
Considérese $-\nabla\cdot(\kappa\nabla u)=f$, con $\kappa>0$, valor nulo en $\Gamma_D$ y $\kappa\nabla u\cdot n=g_N$ en $\Gamma_N$. Multiplicar por una función de prueba $v$ nula en $\Gamma_D$ e integrar por partes conduce a
\begin{equation}
\underbrace{\int_\Omega\kappa\nabla u\cdot\nabla v\,dx}_{a(u,v)}
=\underbrace{\int_\Omega fv\,dx+\int_{\Gamma_N}g_Nv\,ds}_{\ell(v)}.
\label{eq:fem-debil}
\end{equation}
La derivada de segundo orden desaparece del integrando. El flujo natural permanece como contribución de frontera. Para Dirichlet no homogéneo, un levantamiento satisface el dato y la incógnita restante pertenece al espacio homogéneo.

Una malla $\mathcal T_h$ y funciones $\phi_j$ definen $V_h\subset H^1$. Al sustituir $u_h=\sum_j U_j\phi_j$ y elegir sucesivamente $v=\phi_i$,
\begin{equation}
KU=F,\quad K_{ij}=\int_\Omega\kappa\nabla\phi_j\cdot\nabla\phi_i\,dx,
\quad F_i=\ell(\phi_i).
\label{eq:fem-sistema}
\end{equation}
En un triángulo $P_1$, cada función es afín, vale uno en un vértice y cero en los otros. Sus gradientes son constantes dentro del elemento. El ensamblaje suma matrices locales según los nodos compartidos. La estructura dispersa nace del soporte local: dos bases sin elementos comunes no contribuyen a la misma integral \cite{fenics2016}.

\begin{figure}[htbp]\centering
\includegraphics[width=\linewidth]{revision2026/fem_poisson.pdf}
\caption[Recreación de Poisson con FEM triangular.]{Recreación computacional propia de la formulación de Poisson con FEM $P_1$ presentada en el tutorial de FEniCS \cite{fenics2016}, usando $u=\sin(\pi x)\sin(\pi y)$. Se muestran malla, solución calculada y error. El ensamblaje se implementó con SciPy; no se ejecutó FEniCS ni el mecanismo propuesto.}\label{fig:fem}
\end{figure}

Cuando el problema es coercivo, la ortogonalidad de Galerkin relaciona el error con la mejor aproximación disponible en $V_h$. Para solución regular y elementos lineales, es habitual obtener orden uno en seminorma $H^1$ y orden dos en $L^2$ bajo hipótesis adicionales. Singularidades geométricas o interfaces pueden reducir esos órdenes. El error algebraico del sistema resuelto es otro componente y puede ser mucho menor que el error del campo.

\begin{figure}[htbp]\centering
\includegraphics[width=.72\linewidth]{revision2026/fem_convergencia.pdf}
\caption[Convergencia de la recreación FEM.]{Verificación numérica propia de la recreación FEM de la Figura~\ref{fig:fem}, siguiendo el análisis de refinamiento del tutorial \cite{fenics2016}. Tres mallas producen órdenes observados próximos a dos en $L^2$ y uno en $H^1$. Son resultados del ejemplo clásico, no de PIKAN.}\label{fig:fem-orden}
\end{figure}

\subsection{Estimadores a posteriori y refinamiento hp}
Un estimador a posteriori utiliza la solución calculada para localizar incertidumbre de discretización. En Poisson con coeficiente unitario, un indicador residual interior puede escribirse
\begin{equation}
\eta_K^2=h_K^2\|f+\Delta u_h\|_{L^2(K)}^2+
\frac12\sum_{e\subset\partial K\cap\Omega}h_e\|[\nabla u_h\cdot n_e]\|_{L^2(e)}^2.
\label{eq:estimador}
\end{equation}
El primer término mide incumplimiento dentro del elemento; el segundo, incompatibilidad de flujos entre vecinos. La frontera de Neumann aporta su propia discrepancia. La notación de salto utiliza una orientación coherente para comparar trazas. Las cotas de confiabilidad y eficiencia requieren hipótesis; este indicador no es una identidad con el error.

La marcación de Dörfler selecciona elementos $\mathcal M$ que satisfacen
\begin{equation}
\sum_{K\in\mathcal M}\eta_K^2\geq\vartheta\sum_{K\in\mathcal T_h}\eta_K^2,
\qquad 0<\vartheta<1.
\end{equation}
Así concentra una fracción controlada de la señal estimada \cite{dorfler1996}. Refinar $h$ reduce tamaño; refinar $p$ eleva el grado; adaptar $hp$ combina ambas operaciones. Una malla refinada introduce grados de libertad en el espacio físico. En una red, la ubicación de una base puede estar en una coordenada interna, por lo que la analogía con FEM no transfiere automáticamente localidad ni garantías de convergencia.
